\documentclass[12pt]{article}
\usepackage[margin=1in]{geometry}
\usepackage{hyperref}
\pagestyle{empty}
\begin{document}

\begin{flushright}
Huang Zhongchang\\
Independent Researcher\\
Nanjing, China\\
\today
\end{flushright}

\vspace{0.3in}

\noindent The Editor\\
Physical Review D\\
American Physical Society

\vspace{0.2in}

\noindent \textbf{Re: Submission to Physical Review D}

\vspace{0.1in}

\noindent Dear Editor,

I am submitting ``An Information-Capacity Model of Gravity and Quantum Decoherence'' to Physical Review D.

The paper constructs a model in which Newtonian gravity and quantum decoherence emerge as regimes of a single dynamical equation. The model has one irreducible definition---a causal event is an asymmetric influence producing a determinate change (Definition~C)---plus a small set of motivated structural choices. The 1-bit quantization follows from a structural argument (flagged as "not a formal theorem" in the main text). The free-capacity fraction $q$ obeys a unified telegraph equation. In the static limit $\nabla^2(\ln 1/q)=0$, the $1/r$ functional form of the Newtonian potential is derived; the coefficient $GM/c^2$ is calibrated to the observed value of $G$. A postulated metric (causal locking hypothesis) yields PPN parameters $\beta=\gamma=1$, consistent with solar system tests. In vacuum ($q\to1$), state-dependent damping vanishes, restoring Lorentz invariance. An H-theorem is proved via a Lyapunov functional. The decomposition $\nabla^2 q/q=|\nabla q|^2/q^2$ unifies diffusion and self-organization as two regimes of one process.

The quantum-mechanical foundation builds on the peer-reviewed QM reconstruction literature (Zeilinger 1999, Daki\'c \& Brukner 2011, Chiribella/D'Ariano/Perinotti 2011, Masanes \& M\"uller 2011, Shaghaghi 2023) and Zilly (2026). The classical gravitational derivation (Sections~3--7) depends only on Definition~C and structural choices, independently of the QM foundation.

Nine limitations are stated explicitly in the abstract and enumerated in Section~IX: (1)~the damping form is a modeling ansatz, not uniquely derived; (2)~$G$ is calibrated, not predicted; (3)~the spacetime metric is a structural postulate, not derived from Definition~C; (4)~only the Newtonian limit of GR is recovered; (5)~matter coupling is not microscopically derived; (6)~the 1-bit quantization argument is structural, not a formal theorem; (7)~GW170817 consistency requires an unproven tensor-scalar decoupling conjecture; (8)~the volume-scaling entropy bound conflicts with the holographic area-scaling bound; (9)~the continuum limit and $d=3$ are empirical inputs. A tiered-assumption classification (Table~I) is applied to both this framework and Bassani \& Magueijo (PRD 111, 103529, 2025), providing a transparent accounting of all inputs.

I suggest referees with expertise in emergent gravity, quantum foundations, and causal set theory. The manuscript is not under consideration elsewhere.

\vspace{0.15in}
\noindent Respectfully,

\vspace{0.3in}
\noindent Huang Zhongchang

\end{document}
